\devoirDS{5}{Chapitres 14 à 18}{Son et distance/6,Spectre et réfraction/7,Une image nette/7}
\begin{blocLivre}{ds-5-1}\begin{exoBac}{6}{Son et distance}
Deux microphones séparés de 1,02 m reçoivent le même son avec un retard de 3,00 ms. Un enregistrement distinct montre cinq périodes en 10,0 ms.
\schema{retard-son}
\begin{questions}
\item Calculer la vitesse de propagation.
\item Calculer période et fréquence du signal périodique.
\item Expliquer la différence entre hauteur, timbre et niveau sonore.
\end{questions}
\end{exoBac}\end{blocLivre}
\begin{blocLivre}{ds-5-2}\begin{exoBac}{7}{Spectre et réfraction}
Une source possède trois raies à 450, 520 et 650 nm. Le faisceau arrive de l’air sur un verre d’indice 1,50 avec $i=30,0^\circ$.
\schema{refraction}
\begin{questions}
\item La source est-elle monochromatique ? Convertir 520 nm en mètres.
\item Calculer l’angle de réfraction.
\item Quel est l’angle de réflexion ?
\item Pourquoi un prisme peut-il séparer les radiations de la source ?
\end{questions}
\end{exoBac}\end{blocLivre}
\begin{blocLivre}{ds-5-3}\begin{exoBac}{7}{Une image nette}
Une lentille de focale 10 cm reçoit la lumière d’un objet de 2,0 cm situé à 30 cm avant elle.
\schema{lentille-3f}
\begin{questions}
\item Construire l’image avec deux rayons.
\item Lire position et taille de l’image puis calculer le grandissement.
\item Où placer l’écran pour obtenir la netteté ?
\item Associer les éléments du montage au modèle de l’œil réduit.
\end{questions}
\end{exoBac}\end{blocLivre}

